\chapter{Boundary actions and their comparison homotopies}
\label{ch:vol2-boundary-actions}

A bulk observable acts on boundary conditions and on the operators
between them. Its action may vanish on each boundary while remaining
nonzero on the category of boundaries. Two cubic boundaries exhibit
this phenomenon. They lead to a complex that retains the actions,
the homotopies comparing them, and the equations imposed by composition.
The resulting ordinary derived centre gives an algebraic model of
coherent actions. In a mixed holomorphic--topological field theory,
which of these actions come from local bulk observables?

\section{A relation detected between two boundaries}
\label{sec:vol2-cubic-boundaries}

Let $k$ be a field of characteristic zero and put $R=k[x]$.
A \emph{matrix factorization} of $W\in R$ is a finite free
parity-graded $R$-module $E=E^{\bar0}\oplus E^{\bar1}$ with an odd
endomorphism $Q_E$ satisfying $Q_E^2=W\id_E$.

For two factorizations of the same potential, set
\begin{equation}\label{eq:vol2-mf-differential}
 \delta f=Q_Ff-(-1)^{|f|}fQ_E,
 \qquad f\in\Hom_R(E,F).
\end{equation}
Here $|f|\in\Z/2$ is its parity. Since $W$ is central,
$\delta^2f=Wf-fW=0$. Composition satisfies the graded Leibniz rule.
Thus the factorizations form a differential category with parity grading.

For the cubic potential, take the factorizations on $R^{1|1}$
\begin{equation}\label{eq:vol2-cubic-pair}
 Q_X=\begin{pmatrix}0&x\\x^2&0\end{pmatrix},
 \qquad
 Q_Y=\begin{pmatrix}0&x^2\\x&0\end{pmatrix},
 \qquad
 f=\begin{pmatrix}x&0\\0&1\end{pmatrix}:X\longrightarrow Y.
\end{equation}
The map $f$ is even and closed. Multiplication by $x$ has an odd
primitive on each object: there are $h_X,h_Y$ with
$\delta h_X=x\id_X$ and $\delta h_Y=x\id_Y$.
Their compatibility along $f$ is measured by
\[
 c(f)=h_Yf-fh_X.
\]
It is closed, because $\delta c(f)=xf-fx=0$.

\begin{proposition}\label{prop:vol2-cubic-defect}
For the objects and map in \eqref{eq:vol2-cubic-pair}, every choice
of $h_X,h_Y$ gives the same nonzero class
\[
 [c(f)]=[J]\in H^{\bar1}\Hom_R(X,Y),
 \qquad J=\begin{pmatrix}0&1\\-1&0\end{pmatrix}.
\]
In particular, the objectwise primitives cannot be extended to a
nullhomotopy compatible with this morphism.
\end{proposition}

\begin{proof}
The most general primitives are
\[
 h_X=\begin{pmatrix}0&b\\1-xb&0\end{pmatrix},
 \qquad
 h_Y=\begin{pmatrix}0&1-xc\\c&0\end{pmatrix},
 \qquad b,c\in R.
\]
Indeed, an odd matrix with upper entry $u$ and lower entry $v$
is a primitive on $X$ precisely when $xv+x^2u=x$, so
$v=1-xu$. The equation on $Y$ is analogous.

Multiplication gives
\begin{equation}\label{eq:vol2-cubic-defect}
 c(f)=(1-x(b+c))J.
\end{equation}
For an even morphism $a=\operatorname{diag}(a_0,a_1)$ and an odd
morphism with entries $b_0,c_0$, the two differentials are
\[
 \delta a=x(xa_1-a_0)J,
 \qquad
 \delta\begin{pmatrix}0&b_0\\c_0&0\end{pmatrix}
       =x(b_0+c_0)f.
\]
Consequently the even cycles and boundaries are $Rf$ and $xRf$,
and the odd cycles and boundaries are $RJ$ and $xRJ$.
The coefficient of $J$ in \eqref{eq:vol2-cubic-defect} is one
modulo $x$, independently of $b,c$.
\end{proof}

For three objects with such primitives and closed even maps
$f:X\to Y$, $g:Y\to Z$, the same definition gives
\[
 c(gf)=c(g)f+gc(f).
\]
Thus even a primitive for each $c(f)$ would leave a compatibility
equation for composition. The complex of coherent actions must contain
components on morphisms and on longer composable strings.

\section{The ordinary derived centre}
\label{sec:vol2-ordinary-centre}

Let $\mathcal E$ be a small unital differential graded category
over $k$. Complexes have cohomological grading,
$V[1]^j=V^{j+1}$. The same constructions apply with degrees reduced
modulo two to the preceding category of matrix factorizations.
All tensors in this section are over $k$, and total Hom complexes
use products. In particular, no boundedness in the cohomological
degrees is implicit.
Write $D(k)$ for the derived category of coefficient complexes.

The \emph{diagonal bimodule} is
$\Delta_{\mathcal E}(a,b)=\mathcal E(a,b)$, with the two actions
given by composition. Write
$\mathcal E^e=\mathcal E^{\mathrm{op}}\otimes_k^{\mathbb L}\mathcal E$
for the derived enveloping category. The
\emph{ordinary derived centre} and its cohomology are
\begin{equation}\label{eq:vol2-ordinary-centre}
 Z_k(\mathcal E)=
 \RHom_{\mathcal E^e}(\Delta_{\mathcal E},\Delta_{\mathcal E}),
 \qquad
 \HH_k^*(\mathcal E)=H^*Z_k(\mathcal E).
\end{equation}
This definition fixes the coefficient tensor. Relative centres over
a larger ring will use that ring's tensor and derived Hom instead.

To compute \eqref{eq:vol2-ordinary-centre}, resolve the diagonal
by its two-sided bar. After bimodule Hom, its cochains have factors
\begin{equation}\label{eq:vol2-all-object-bar}
 \prod_{n\ge0}\prod_{w_0,\ldots,w_n}
 \Hom_k\!\left(
 \mathcal E(w_{n-1},w_n)[1]\otimes\cdots\otimes
 \mathcal E(w_0,w_1)[1],\mathcal E(w_0,w_n)\right).
\end{equation}
At $n=0$ the input is $k$. Every intermediate object is an index
of the product. The normalized cochains vanish when an input is an
identity. Denote their product totalization by
$\CC_k^\bullet(\mathcal E)$.

The differential has an internal part and a part that composes
adjacent inputs or acts at an endpoint. An unsuspended convention
makes their signs explicit. An arity-$n$ cochain $F$ of internal
degree $r$ has total degree $r+n$. Put $\partial F=[\delta,F]$,
the graded Hom differential.

For homogeneous composable inputs,
ordered so that the rightmost acts first, set
\begin{align}\label{eq:vol2-hochschild-signs}
 (b_HF)(a_1,\ldots,a_{n+1})
 &=(-1)^{r|a_1|}a_1F(a_2,\ldots,a_{n+1})\notag\\
 &\quad+\sum_{j=1}^n(-1)^j
 F(a_1,\ldots,a_ja_{j+1},\ldots,a_{n+1})\notag\\
 &\quad+(-1)^{n+1}F(a_1,\ldots,a_n)a_{n+1},\\
 (\mathcal DF)_n&=(-1)^n\partial F_n+b_HF_{n-1},
 \qquad F_{-1}=0.\label{eq:vol2-hochschild-total}
\end{align}
Associativity and the Leibniz rule give
$b_H^2=\partial^2=0$ and $\partial b_H=b_H\partial$.
Hence $\mathcal D^2=0$. Each component involves only two arities,
so the differential is defined on the product totalization.

\begin{proposition}[Bar representation and evaluation]
\label{prop:vol2-centre-evaluation}
The normalized complex $\CC_k^\bullet(\mathcal E)$ represents
$Z_k(\mathcal E)$. For every coefficient complex $P\in D(k)$,
evaluation induces a natural equivalence
\begin{equation}\label{eq:vol2-centre-evaluation}
 \RHom_k(P,Z_k(\mathcal E))
 \simeq
 \RHom_{\mathcal E^e}
 (P\otimes_k^{\mathbb L}\Delta_{\mathcal E},\Delta_{\mathcal E}).
\end{equation}
\end{proposition}

\begin{proof}
Use cofibrant models for the coefficient complexes. The augmented
two-sided bar contracts by insertion of a unit at an outside end.

Its filtration by word length has induced bimodules as successive
quotients. Over the field $k$, their coefficient complexes are
homotopically projective: Hom from them into an acyclic complex
is acyclic. A finite length stage therefore has the same property.

Hom out of the full direct-sum bar is the inverse limit of the
Hom complexes of these stages. The transition maps are surjective
in every degree, since the length filtration splits on graded
modules. The kernels are acyclic by the same coefficient argument,
so primitives can be lifted successively through the tower.
Its inverse limit is therefore acyclic. Thus the bar is a
homotopically projective resolution.

The degeneracy contraction splits off the cochains containing an
identity, leaving the normalized complex. Applying bimodule Hom
to the remaining resolution gives \eqref{eq:vol2-all-object-bar}.
For a projective model of $P$, the tensor--Hom adjunction on each
bar term identifies the two sides of
\eqref{eq:vol2-centre-evaluation}. Products commute with this
adjunction, and the identifications commute with the differential.
\end{proof}

The identity action is the nullary family $\id_w$.
The cup product composes outputs after cutting a composable word.
The brace operations insert the output of one cochain into an input
of another, with the Koszul signs of the suspended inputs.
These are operations on the full cochain complex. Retaining only
the cohomology of the endomorphisms of each object removes the
positive-arity terms where Proposition~\ref{prop:vol2-cubic-defect}
found its obstruction.

\section{The scalar action on all cubic boundaries}
\label{sec:vol2-cubic-scalar-centre}

The cubic calculation extends to every finite free boundary.
It distinguishes a relation in the category from a relation on each
of its objects. Put $R=k[[x]]$. Let $\mathcal C$ be the
parity differential category of all matrix factorizations of $x^3$
on standard finite free parity-graded $R$-modules. Fixing standard
modules makes the objects a set. Its morphisms are $R$-linear,
with differential \eqref{eq:vol2-mf-differential}.

Here the centre is \emph{absolute over $k$}: its cochains are
$k$-multilinear, with no continuity condition on those maps.
They are normalized and product-totalized as above. For $g\in R$,
scalar multiplication defines a cocycle $\rho(g)$ whose nullary
value is $g\id_E$ and whose higher values vanish.

\begin{theorem}\label{thm:vol2-cubic-scalar-kernel}
The kernel of the scalar action
\[
 R\longrightarrow\HH_k^{\bar0}(\mathcal C),
 \qquad g\longmapsto[\rho(g)],
\]
is $(x^2)$. Nevertheless, multiplication by $x$ is nullhomotopic
on every object of $\mathcal C$.
\end{theorem}

\begin{proof}
For a cochain primitive with nullary component $h_E$ and unary
component $u$, equations \eqref{eq:vol2-hochschild-signs} and
\eqref{eq:vol2-hochschild-total} begin with
\begin{align}
 \delta h_E&=g\id_E,\label{eq:vol2-scalar-primitive-zero}\\
 \delta u(f)-u(\delta f)
 &=(-1)^{|f|}fh_E-h_Ff.
 \label{eq:vol2-scalar-primitive-one}
\end{align}
The map $h_E$ is odd and $u$ preserves parity.
If all later components vanish, the binary equation is
$u(vf)=u(v)f+vu(f)$.

Differentiate the matrices in their fixed bases and take
\[
 h_E=\frac13Q'_E,
 \qquad u(f)=\frac13f',
 \qquad F_n=0\quad(n\ge2).
\]
Differentiating $Q_E^2=x^3\id_E$ proves
\eqref{eq:vol2-scalar-primitive-zero} for $g=x^2$.
For a homogeneous morphism,
\[
 (\delta f)'=\delta(f')+Q'_Ff-(-1)^{|f|}fQ'_E,
\]
which proves \eqref{eq:vol2-scalar-primitive-one}.
The derivative kills identities and satisfies the Leibniz rule,
so every remaining component of $\mathcal DF$ vanishes.
Multiplying the components by any $a\in R$ gives a primitive
for $ax^2$, since both cochain differentials commute with the
central scalar action on the output.

Conversely, reduce $g$ modulo $(x^2)$ to $\alpha+\beta x$.
On the object $X$ of \eqref{eq:vol2-cubic-pair}, every odd
matrix has the form
\[
 h=\begin{pmatrix}0&b\\d&0\end{pmatrix},
 \qquad \delta h=x(d+xb)\id_X.
\]
The nullary equation forces $\alpha=0$ and $d+xb=\beta$.

The odd endomorphism
\[
 \psi=\begin{pmatrix}0&1\\-x&0\end{pmatrix}
\]
is closed. The unary equation on $\psi$ would give
\[
 \delta u(\psi)=-(\psi h+h\psi)=-(\beta-2xb)\id_X.
\]
Every matrix boundary on $X$ is divisible by $x$.
Evaluation at $x=0$ forces $\beta=0$. Higher arities do not
enter these two equations, so they cannot remove the obstruction.

It remains to prove the objectwise assertion for arbitrary rank.
Over the fraction field of $R$, the two odd blocks of $Q_E$
are inverse up to $x^3$. The even and odd ranks therefore agree.
Smith normal form over the discrete valuation ring $R$
diagonalizes one block. The equation $Q_E^2=x^3\id_E$ forces
the other block to be diagonal in the complementary bases.
The factorization is consequently a direct sum of blocks
$(x^a,x^{3-a})$, where $0\le a\le3$.

For $a=0,3$ one block is invertible, so the object is contractible.
For $a=1$ the matrix
\[
 H=\begin{pmatrix}0&0\\1&0\end{pmatrix}
 \qquad\text{satisfies}\qquad \delta H=x\id.
\]
Exchanging the two parity summands gives a primitive for $a=2$.
Direct sums and changes of basis preserve these homotopies.
\end{proof}

The scalar image in this absolute centre is therefore
$k[[x]]/(x^2)$. This computes the scalar image, without asserting
that it is the entire centre. The coefficient convention enters
the proof at the unary map:
\[
 u(xf)=\frac13f+xu(f).
\]
It is $k$-linear and is not $R$-linear. A relative centre over
$R$ asks a different question, since it excludes this primitive.

\section{Two scalar endpoints joined by a complex}
\label{sec:vol2-triangular-centre}

The cubic example required an infinite bar of a category.
There is also a finite model in which the entire relative centre,
including its multiplication in a family, can be computed.

Let $R$ be a commutative ring and $M$ a bounded complex of finite
projective $R$-modules. Use integer cohomological degrees and set
\begin{equation}\label{eq:vol2-triangular-algebra}
 T(M)=Re_0\oplus Re_1\oplus M,
 \qquad e_1Me_0=M,
 \qquad M^2=0.
\end{equation}

The idempotents are orthogonal, closed, and sum to the unit.
All other idempotent actions on $M$ vanish. The differential on
$M$ is its given differential. Thus $M$ is the complex of arrows
from object zero to object one, and the endpoint endomorphisms
are scalar. All tensors in this section are over $R$ unless
another ring is displayed.

Put $T=T(M)$ and $T^e=T\otimes_R T^{\mathrm{op}}$, with graded
opposite multiplication. The stated projectivity makes this a
model of the derived enveloping algebra. Define
\[
 Z_R(T)=\RHom_{T^e}(T,T),
 \qquad
 E=\End_R(M),
 \qquad d_EF=d_MF-(-1)^{|F|}Fd_M.
\]
The difference of the two scalar actions is the chain map
\[
 \chi:R^2\longrightarrow E,
 \qquad \chi(a,b)=(a-b)\id_M,
\]
where $a$ acts at the source and $b$ at the target.
Write $D(R)$ for the derived category of $R$-module complexes.

\begin{theorem}\label{thm:vol2-triangular-cone}
The relative centre is represented by the homotopy fibre
\begin{equation}\label{eq:vol2-triangular-cone}
 Z_R(T(M))\simeq\operatorname{Cone}(\chi)[-1].
\end{equation}
With the fibre sign convention, its terms and differential are
\[
 Z^n=(R^2)^n\oplus E^{n-1},
 \qquad d_Z(a,b,F)=(da,db,\chi(a,b)-d_EF).
\]
For every $P\in D(R)$, evaluation gives
\begin{equation}\label{eq:vol2-relative-evaluation}
 \RHom_R(P,Z_R(T))
 \simeq\RHom_{T^e}(P\otimes_R^{\mathbb L}T,T).
\end{equation}
\end{theorem}

\begin{proof}
Let $S=Re_0\oplus Re_1$. The normalized relative bimodule bar
of $T$ over $S$ stops at length one, since $M\otimes_SM=0$:
\begin{equation}\label{eq:vol2-short-bar}
 T\otimes_SM\otimes_ST
\longrightarrow T\otimes_ST\longrightarrow T.
\end{equation}

The first term has the extra degree $-1$ of the bar suspension.
Choose the boundary on an arrow to be
$e_1\otimes m-m\otimes e_0$, with the corresponding suspended
internal differential. Multiplication adds the two copies of
$M$ in the middle term. Their difference is its kernel.
On the idempotents, multiplication is the identity.
This proves exactness of the augmented sequence of complexes.

The terms are induced from the split idempotent bimodules of
$S$, with bounded projective coefficients $R$ or $M$.
They are therefore homotopically projective as $T$-bimodules.
Applying bimodule Hom into $T$ yields $R^2$ at length zero
and $\End_R(M)$ at length one. The two endpoint actions give
$(a-b)\id_M$, and suspension gives $-d_E$.
This is \eqref{eq:vol2-triangular-cone}.
Tensoring the same bar with a projective model of $P$ and using
tensor--Hom adjunction proves \eqref{eq:vol2-relative-evaluation}.
\end{proof}

\begin{example}[Parallel arrows]\label{ex:vol2-parallel-arrows}
For $R=k$ and $M=k^m$ in degree zero, with $m>0$, the centre
complex is
\[
 k^2\longrightarrow\operatorname{Mat}_m(k),
 \qquad (a,b)\longmapsto(a-b)\id,
\]
in degrees zero and one. Hence
\[
 \HH_k^0(T)=k,
 \qquad \HH_k^1(T)=\operatorname{Mat}_m(k)/k\id,
 \qquad \HH_k^j(T)=0\quad(j\ne0,1).
\]
Insertion of one unary cochain into another is composition on
$M$, so the degree-one Gerstenhaber bracket is the matrix
commutator modulo scalars. The product of two positive-degree
classes vanishes: a cup product would require two composable
arrows, and there is no such word. Thus a single arrow has no
positive centre cohomology, whereas three parallel arrows give
eight independent degree-one classes.
\end{example}

\section{A differential varying with a parameter}
\label{sec:vol2-varying-centre}

The joining complex can change while both endpoint algebras
remain fixed. Put $R=k[t]$ and take
\[
 M_t=[Ru\xrightarrow{\ t\ }Rv],
 \qquad |u|=-1,
 \qquad |v|=0.
\]
Its endomorphism complex is
\[
 E^{-1}=R\xrightarrow{s\mapsto(ts,ts)}E^0=R^2
 \xrightarrow{(p,q)\mapsto t(p-q)}E^1=R.
\]
Here $s$ sends $v$ to $su$, the pair $(p,q)$ acts on $(u,v)$,
and the last scalar sends $u$ to a multiple of $v$.
The centre in Theorem~\ref{thm:vol2-triangular-cone} is therefore
\begin{equation}\label{eq:vol2-varying-complex}
 R^3\xrightarrow{D_0}R^2\xrightarrow{D_1}R,
 \qquad
 \begin{aligned}
 D_0(a,b,s)&=(a-b-ts,a-b-ts),\\
 D_1(p,q)&=-t(p-q),
 \end{aligned}
\end{equation}
in degrees zero, one, and two.

For a complex of free $R$-modules, its \emph{Bockstein}
on cohomology modulo $t$ sends $[\bar c]$ to
$[t^{-1}dc\bmod t]$, where $c$ lifts the cocycle $\bar c$.
This is the connecting map for the first-order thickening
$k[t]/(t^2)$, with the submodule $(t)$ identified with $k$.

\begin{theorem}\label{thm:vol2-varying-centre}
The full relative Hochschild cohomology algebra is
\begin{equation}\label{eq:vol2-varying-ring}
 \HH_R^*(T(M_t))
 \cong\frac{R[z,\omega]}
 {(z^2-tz,\ t\omega,\ z\omega,\ \omega^2)},
 \qquad |z|=0,
 \qquad |\omega|=2.
\end{equation}
After inverting $t$, it is $R[t^{-1}]\times R[t^{-1}]$ in
degree zero. The derived fibre
$Z_R(T(M_t))\otimes_R^{\mathbb L}k$ has degree-zero algebra
$k[z]/(z^2)$, together with one class $\epsilon$ in degree one
and one class $\omega$ in degree two. In that fibre,
\[
 z\epsilon=z\omega=\epsilon^2=\epsilon\omega=\omega^2=0.
\]
Its Bockstein is
\begin{equation}\label{eq:vol2-centre-bockstein}
 \beta(\epsilon)=-\omega,
 \qquad \beta(z)=0.
\end{equation}
\end{theorem}

\begin{proof}
A degree-zero cycle obeys $a-b=ts$, so its free coordinates
are $(b,s)$. There are no incoming boundaries.
The two cuts of a one-arrow word give the cup product
\begin{equation}\label{eq:vol2-centre-cup}
 (a,b,s)(a',b',s')=(aa',bb',bs'+sa').
\end{equation}
The unit is $(1,1,0)$, and $z=(t,0,1)$ satisfies $z^2=tz$.
Every degree-zero cycle is uniquely $b\,1+s\,z$.

Since $t$ is a non-zero-divisor, $\ker D_1$ is the diagonal
copy of $R$, which equals $\im D_0$. Thus $H^1Z=0$.
The last cokernel is $R/(t)$, generated by the one-arrow
cochain $\omega:u\mapsto v$. Products of two one-arrow
cochains vanish. The left product $z\omega$ uses the target
scalar zero, while the right product uses the source scalar
$t$ and vanishes in $R/(t)$. These are all relations in
\eqref{eq:vol2-varying-ring}, by the computed cohomology groups.

Every term of \eqref{eq:vol2-varying-complex} is free over $R$,
so reduction modulo $t$ computes the derived fibre.
The reduced differential is
\[
 D_0(a,b,s)=(a-b,a-b),
 \qquad D_1=0.
\]
The cohomology dimensions are $2,1,1$ in degrees $0,1,2$.
The degree-one class is represented by $(1,0)\in Z^1$.
At $t=0$, the representative of $z$ has zero endpoint components,
so its products with $\epsilon$ and $\omega$ contain two
one-arrow cochains and vanish. Formula~\eqref{eq:vol2-centre-cup}
also gives $z^2=0$.

The lift $(1,0)$ of $\epsilon$ has differential $-t$.
The lift $(t,0,1)$ of $z$ is already closed.
This proves \eqref{eq:vol2-centre-bockstein}.
After inverting $t$, the roots $0,t$ of $z(z-t)$ separate
and the torsion class disappears.
\end{proof}

The map $b\,1+s\,z\mapsto(b+ts,b)$ identifies the degree-zero
centre with
\[
 R\times_{R/(t)}R.
\]
Its spectrum consists of two affine lines joined at the origin.
The degree-two torsion has a further effect: derived specialization
produces $\epsilon$, which tensoring the cohomology groups alone
with $k$ would omit. At the special fibre, $z$ itself is a nonzero
central class with zero endpoint components. It is carried entirely
by the homotopy on the joining complex.

\section{Gluing endpoints with their own operations}
\label{sec:vol2-dg-triangular-centre}

For differential graded endpoint algebras, the derived centres
act on the joining bimodule, and their difference gives the same
compatibility problem. Work over the field $k$. Let $A_0,A_1$
be unital differential graded algebras and let $M$ be an
$A_1$--$A_0$ bimodule. Form the two-object category
$\mathcal E_M$ with these endomorphism algebras,
$\mathcal E_M(0,1)=M$, and $\mathcal E_M(1,0)=0$.
Use derived coefficient models throughout.

Tensoring an endomorphism of the diagonal $A_0$-bimodule with
$M$ gives its right characteristic action on $M$.
The left action is defined similarly using $A_1$.
Write $\chi$ for right action minus left action.

\begin{proposition}\label{prop:vol2-dg-triangular}
There is a fibre triangle of complexes
\begin{equation}\label{eq:vol2-dg-triangular}
 \CC_k^\bullet(\mathcal E_M)\longrightarrow
 \CC_k^\bullet(A_0)\oplus\CC_k^\bullet(A_1)
 \xrightarrow{\ \chi\ }
 \RHom_{A_1\otimes_k^{\mathbb L}A_0^{\mathrm{op}}}(M,M).
\end{equation}
The cochains use product totalization over all bar lengths.
\end{proposition}

\begin{proof}
Normalize the two-sided bar relative to the object idempotents.
A composable word stays in one corner or crosses once from
zero to one. Put $\overline A_i=A_i/k1$.

The crossing cochains have factors
\[
 \Hom_k\!\left(
 (\overline A_1[1])^{\otimes p}\otimes M[1]\otimes
 (\overline A_0[1])^{\otimes q},M\right),
 \qquad p,q\ge0.
\]
The double bar resolving $M$ as an $A_1$--$A_0$ bimodule
is the direct sum of
\[
 A_1\otimes(\overline A_1[1])^{\otimes p}\otimes M
\otimes(\overline A_0[1])^{\otimes q}\otimes A_0.
\]

Its differential combines the internal differentials,
adjacent products, and the two actions on $M$.
Unit insertion contracts the augmented bar, first on one side
and then on the other. Its length filtration has induced
bimodules with homotopically projective coefficient complexes
as successive quotients. The surjective inverse-tower argument
of Proposition~\ref{prop:vol2-centre-evaluation} proves that
the whole double bar is homotopically projective, including
when the coefficient complexes are unbounded.

Hom from this resolution into $M$ computes the right-hand term
of \eqref{eq:vol2-dg-triangular}. The crossing factor has one
additional suspension, so the kernel of restriction to the
corners is this complex shifted by $-1$.
The remaining endpoint mergers give the right and left
characteristic actions with opposite signs. The resulting
short exact sequence splits on graded cochains and gives the
stated fibre triangle.
\end{proof}

This is a triangle of complexes. The cup products and braces
remain on the Hochschild cochains and their composable strings.
The last term can be replaced by an ordinary endomorphism
complex only when the chosen bimodule model computes derived Hom.
Likewise, selecting fewer boundary objects computes the centre
of that selected category. Comparison with all boundaries is
a further assertion.

\section{From coefficient actions to local bulk observables}
\label{sec:vol2-centre-geometric-comparison}

The bar evaluation gives a precise universal property over the
declared coefficient ring. In a mixed holomorphic--topological theory,
operators vary on a curve and compose in an ordered topological
direction. Their products may acquire terms supported on collision
divisors. A comparison must preserve these operations and their
coefficient modules as well as the ordinary composition of morphisms.

The physical source must be defined independently. Its fields,
action, boundary conditions, and quantization determine a complex
of bulk observables and their action on boundary insertions.
Identifying that complex with a geometric centre requires a geometric
category of boundary probes and its coefficient tensor.
The evaluation map must be compatible with the supported operations.
The comparison must also retain homotopies on morphisms and their
compatibility with composition, as the cubic example requires.

The ordinary centres above supply complete algebraic models of that
compatibility, with evaluation maps and explicit dependence on
coefficients. Their identification with a native chiral centre
remains a separate construction. A geometric issue is already visible
before quantization. Coefficients with poles along a divisor and
coefficients supported on it do not form one ordinary algebra of
functions. Their respective modules and the operations between them
must be specified.
